\makeatletter\def\input@path{{../shared/}{../diagrams/}}\makeatother
\documentclass[11pt,a4paper]{article}
\usepackage{spirdoc}
\SetDoc{SPIR-REL-001}{Release Notes and Revision History}{1.0}{Issued for review}
\Revisions{%
  1.0 & 27 Nov 2026 & \SPIRAuthor & First issue, covering releases 0.1.0 to \SPIRRelease.\\
}
\newcommand{\release}[3]{\subsection*{Release #1\hfill{\normalsize\mdseries\color{spirink2}#2}}\addcontentsline{toc}{subsection}{Release #1}{\small\color{spirink2}Commits: #3}\par}
\newcommand{\commit}[1]{\code{#1}}

\begin{document}
\MakeCover{What changed in each release of SPIR, why it changed, and what it means for users and for the next phase.}
\DocumentControl
\tableofcontents
\clearpage

\section{Overview}
SPIR was built in five releases between 25 and 26 November 2026, each answering feedback from the previous one.
Version numbers follow semantic versioning (\code{major.minor.patch}); the commit IDs refer to the
\href{\SPIRRepoURL}{\SPIRRepo} repository, branch \code{main}.

\Diagram[0.3]{release-timeline}{Release timeline.}

{\small\zebra
\begin{tabularx}{\linewidth}{@{}L{1.6cm}L{2.2cm}Y@{}}
  \thead\hd{Release} & \hd{Date} & \hd{Theme}\\
  0.1.0 & 25 Nov 2026 & Complete proof of concept from spec v1.1: registry, approval workflow, QR resolution, label OCR check, demo roles.\\
  0.1.1 & 26 Nov 2026 & Codespaces start-up fix.\\
  0.2.0 & 26 Nov 2026 & Parts identified the way real receiving labels do; open mode with CSV import, QR labels and scanning.\\
  0.3.0 & 26 Nov 2026 & QR labels that phone cameras read reliably; company logo; A4 sticker sheets.\\
  0.4.0 & 26 Nov 2026 & Approved cozy-glass visual design and brand palette.\\
  0.4.1 & 26 Nov 2026 & Database set-up helpers.\\
  \textbf{0.5.0} & 26 Nov 2026 & \textbf{Two portals (supplier, company) with sign-in; roles reduced to two; saves live at once; QR scans open to anyone.}\\
\end{tabularx}}

\section{Releases}

% ------------------------------------------------------------------------------------------
\release{0.5.0 --- Two portals}{26 November 2026}{\commit{05c8a8d}}
\paragraph{Why} The stakeholder asked for a supplier page to log in, enter product data and generate QR codes; a
company page to scan QR codes and see the information; and only two roles for now, with role-based access to follow.

\paragraph{What changed}
\begin{itemize}
  \item \textbf{Supplier portal} at \code{/\#/supplier} and \textbf{company portal} at \code{/\#/company}, each with
    its own sign-in. Both use \code{admin} / \code{admin} in this phase.
  \item \textbf{Supplier profile}: the supplier account belongs to one supplier, set once on the Profile page. Every
    shipment and QR label carries that name. A supplier column in a CSV is ignored.
  \item \textbf{Shipment form} to add or edit one shipment, with field-level messages. It goes through the same engine
    as the CSV import, so both follow identical rules. Part number, PO and line identify a shipment and its QR code and
    cannot change in an edit.
  \item \textbf{Saves are live immediately}: no draft or review step. A changed save adds a new version under the same
    QR code; the previous version is kept as SUPERSEDED.
  \item \textbf{Certificates} can be uploaded at any time to the live version of a shipment.
  \item \textbf{QR scans open to anyone}: a phone camera opens a public page with the shipment and large certificate
    buttons, no sign-in. A scanned token only opens its own record's certificates; lookups are rate limited.
  \item \textbf{Company portal}: QR scanner, search across all suppliers (PO, part, heat, record), label photo check and
    its history.
  \item \textbf{Removed}: role switcher, review queue, approval workflow, open mode, products / audit / settings screens.
    The previous design is kept on branch \code{archive/roles-and-review}. Old user accounts stay in the database,
    inactive, because history points at them.
  \item Database migration \code{006\_two\_portals.sql}: roles SUPPLIER and COMPANY, accounts \code{supplier-admin} and
    \code{company-admin}, supplier contact fields.
  \item Sessions are HMAC-signed tokens valid for 12 hours. The web app keeps one token per portal and chooses it from
    the address, so the two portals never mix in one browser (a defect found and fixed during testing).
  \item Tests rewritten around the portals: 63 automated tests pass.
\end{itemize}
\paragraph{Impact} Requirements \req{FR-WF-001}--\req{FR-WF-006} and \req{FR-WF-009} move to \emph{deferred};
\req{FR-QR-003} and \req{FR-DOC-005} become deliberate deviations (see SPIR-RTM-001 and SPIR-SCP-001). In the demo data, a
shipment that had a pending correction now shows the corrected values.

% ------------------------------------------------------------------------------------------
\release{0.4.1 --- Database set-up helpers}{26 November 2026}{\commit{0797b99}, \commit{c0ab908}}
\begin{itemize}
  \item The \code{db:*} scripts explain a missing or placeholder \code{DATABASE\_URL} in plain words instead of failing
    with a network error (\code{getaddrinfo ENOTFOUND HOST}).
  \item \code{npm run db:connect}: paste the Neon connection string at a hidden prompt; it is tested and written to
    \code{.env}. Nothing is written if the connection fails.
  \item The example uses \code{sslmode=verify-full}, which removes the driver's SSL warning.
\end{itemize}

% ------------------------------------------------------------------------------------------
\release{0.4.0 --- Cozy glass visual system}{26 November 2026}{\commit{dba8a2b}}
\paragraph{Why} The stakeholder specified a hybrid glassmorphic design. Mockups were made and reviewed first, then
the colour palette and glass treatment were applied to the application.
\begin{itemize}
  \item Deep-blue glass navigation rail (\code{rgba(0,75,141,.92)}, blur 18\,px) over a warm mist canvas
    (\code{\#F4F6F9}) with faint amber and sky-cyan washes.
  \item Frosted white cards, tables and panels (\code{rgba(255,255,255,.78)}, blur 14\,px, 18\,px corners); brand-blue
    table header rows and primary buttons.
  \item Brand status colours with readable text: APPROVED green, PENDING amber, DRAFT yellow, REJECTED red,
    SUPERSEDED grey. Colours too light for small text colour the dot or border, and a darker shade of the same hue
    carries the words (4.5:1 contrast).
  \item Figtree interface font, IBM Plex Mono for identifiers; warm amber hang tag and brand mark.
  \item Prints on plain white; a fallback for browsers without backdrop blur.
  \item Fixed: a Menu button showing on desktop; misaligned cards on the home page.
\end{itemize}

% ------------------------------------------------------------------------------------------
\release{0.3.0 --- QR labels that phone cameras read}{26 November 2026}{\commit{d312f55}}
\paragraph{Why} Users reported that phone cameras could not detect the printed QR codes, and asked for codes
personalised for the company.
\paragraph{Root causes found} images scaled down with pixel doubling (distorted squares), a 1-module white border
(the standard asks for 4), small codes on screen, and long links (Codespace address plus a 43-character token).
\begin{itemize}
  \item Tokens shortened to 16 characters of Crockford base32 (80 random bits). Older tokens still resolve.
  \item Error correction chosen per code (strongest level that does not add squares); vector SVG output; white border of
    at least 4 modules including the label padding.
  \item Company logo in a grid-aligned cut-out in the centre of the code (\fpath{web/public/brand/logo.png} or
    \fpath{logo.svg}). Every code is decoded again with ZXing after drawing; if the logo gets in the way it is
    shrunk or left out, and the page says so.
  \item A4 sticker sheets: 8 (99.1\,×\,67.7\,mm, Avery L7165), 14 (L7163) or 21 (L7160) per sheet, start position for
    part-used sheets, printer alignment offset, outlines for test prints; 300\,dpi PNG per label.
  \item The label shows Part \#, PO \# with line and Heat \#, as requested.
\end{itemize}
\paragraph{Measured result} At the smallest label size with a Codespace-length link and a blurred, tilted capture,
decoding went from 13 of 20 (old codes) to 20 of 20 (new codes with logo).

% ------------------------------------------------------------------------------------------
\release{0.2.0 --- Real receiving labels; open mode}{26 November 2026}{\commit{8340c74}, \commit{9ee6eac}}
\paragraph{Identification fields (\commit{8340c74})} From photos of real company receiving labels and supplier labels:
\begin{itemize}
  \item Mandatory: part number (company item number), purchase order number and PO line, heat number, quantity,
    pattern number, supplier. Optional: radiographic test number, part description, material, design revision.
  \item \code{02/02} on a label means piece 2 of a quantity of 2; ``Rev:E'' is the design revision; ``LINE 4'' the PO line.
  \item New attribute set and label profile \code{RECEIVING\_LABELS v1} (aliases for both label types, patterns,
    fallbacks); comparison rules \code{compare-v2}, record resolver \code{resolver-v2}.
  \item Crop tool and automatic upscaling for photos; golden test corpus of real Tesseract and ZXing output; a private,
    git-ignored corpus for real photos.
\end{itemize}
\paragraph{Open mode (\commit{9ee6eac})} The stakeholder asked to remove all actors for now: import product data from a
CSV, generate QR codes first, and scan them to retrieve the information.
\begin{itemize}
  \item No sign-in; one generic operator. Import CSV (all-or-nothing, every problem listed by row; a changed row adds a
    version under the same QR code). QR labels page (print sheet, PNG, CSV of links). Scan page.
\end{itemize}

% ------------------------------------------------------------------------------------------
\release{0.1.1 --- Codespaces fix}{26 November 2026}{\commit{dc76ba8}}
The Codespace opened in recovery mode with \code{npm: command not found}. The dev container now uses a Debian base
image plus the official Node 24 feature.

% ------------------------------------------------------------------------------------------
\release{0.1.0 --- Proof of concept from spec v1.1}{25 November 2026}{\commit{8491989}}
\begin{itemize}
  \item Registry: suppliers, products with versioned governed attributes (Postgres JSONB), traceability records with
    heats, batches (many-to-many) and serials, documents in a private Neon Object Storage bucket with 5-minute links.
  \item Approval workflow: DRAFT, PENDING, APPROVED, REJECTED, SUPERSEDED, WITHDRAWN; review queue with differences;
    mandatory rejection reason; one approved and one open version per record (database indexes); approved versions
    immutable and the audit trail append-only (database triggers).
  \item QR tokens: random, stored hashed and encrypted; rotate and revoke; resolution only for signed-in company users
    and only of approved data.
  \item Label OCR check: Tesseract.js and ZXing in the browser; label profiles; extraction, normalization, per-field
    comparison; nine field results and five overall statuses; operator review; discrepancies; camera and
    keyboard-wedge scanners.
  \item Demo role switcher with supplier, reviewer, administrator, QR user, inspector and auditor personas.
  \item Editable seed CSV (valve castings, fictional suppliers), GitHub Codespaces, GitHub Actions (CI, Pages, Neon
    deploy); schema applied to Neon.
\end{itemize}

\section{Upgrade notes}
\begin{itemize}
  \item Run \code{npm run db:setup} (or \code{npm run db:migrate}) after pulling a release; migrations are numbered and
    applied once. The Neon production branch had migration 001 only at the time of writing; 002 to 006 are pending.
  \item Printed QR labels survive every release: tokens are never changed by an upgrade, and old 43-character tokens
    still resolve.
  \item From 0.5.0, \code{AUTH\_MODE} in \code{.env} has no effect and can be removed.
\end{itemize}
\end{document}
